% Edited conversion of source pp. 39--40.
\ReportHeading
  {Дослідження параметрів руйнування для інтерфейсної тріщини в п’єзоелектричному біматеріалі, розташованому у скінченній області}
  {В.~Я. Адлуцький}
  {Дніпровський національний університет імені Олеся Гончара}
  {}

Розглянуто задачу плоскої деформації для двошарової прямокутної області, заповненої п’єзоелектричними матеріалами з різними напрямками поляризації. Шари жорстко зчеплені вздовж інтерфейсу, за винятком центрально розташованої міжфазної тріщини. На верхній і нижній границях області задано рівномірно розподілені розтягувальні зусилля та електричні заряди. Через відмінність властивостей компонентів біматеріалу на бокових границях додатково задають кусково-сталі нормальні зусилля, які забезпечують неперервність деформації вздовж інтерфейсу.

Тріщина вважається електронепроникною та вільною від механічних зусиль. Шуканими параметрами руйнування є розкриття тріщини та швидкість вивільнення енергії. Досліджено вплив скінченних розмірів області на ці параметри та проведено порівняння з точними розв’язками для нескінченної області.

Крайову задачу розв’язано методом скінченних елементів. Швидкість вивільнення енергії визначено інтегральним методом віртуального закриття тріщини з постпроцесорною обробкою чисельних результатів, аналогічною до підходів у~\cite{adlucky:2024jms,adlucky:2024asm}.

Чисельний аналіз показав, що коли характерні розміри скінченної області приблизно на порядок перевищують довжину тріщини, результати добре узгоджуються з аналітичними даними для нескінченної області. Виконання або невиконання умови неперервності деформації на інтерфейсі майже не впливає на величину розкриття тріщини, але істотно змінює швидкість вивільнення енергії. За дотримання неперервності чисельні значення добре збігаються з аналітичними; у разі її порушення виникає додаткова концентрація електромеханічних полів біля вершин тріщини. Встановлено також помітний вплив напрямків поляризації матеріалів на параметри руйнування.

\textit{Роботу виконано за фінансової підтримки Національного фонду досліджень України, грант №~2025.07/0117.}

\begin{thebibliography}{9}
\bibitem{adlucky:2024jms}
V.~J. Adlucky, M.~S. Levchenko, and V.~V. Loboda,
``Finite-element analysis of the parameters of fracture in a piezoelectric bimaterial with interface crack for various types of boundary conditions on its faces,''
\emph{Journal of Mathematical Sciences}, vol.~279, no.~2, pp.~1--16, 2024.

\bibitem{adlucky:2024asm}
V.~Adlucky, N.~Guk, and V.~Loboda,
``Finite element modeling of partially conductive interfacial crack in piezoelectric bimaterial,''
in \emph{Selected Problems of Solid Mechanics and Solving Methods}, Advanced Structured Materials, vol.~204.
Cham: Springer, 2024, pp.~15--30.
\end{thebibliography}
